\documentclass[10pt,a4paper]{article}
\usepackage[T1]{fontenc}
\usepackage[margin=2cm]{geometry}
\usepackage{isabelle,isabellesym}
\usepackage{amsmath,amssymb}
\usepackage{pdfsetup}

\urlstyle{rm}
\isabellestyle{it}

\begin{document}

\title{Kirchhoff's Matrix-Tree Theorem}
\author{Carles Mar\'in}
\maketitle

\begin{abstract}
For a finite simple graph and any vertex $r$, the determinant of the reduced Laplacian, the Laplacian with the row
and the column of $r$ deleted, is the number of spanning trees of the graph. We formalize this theorem of
Kirchhoff~\cite{Kirchhoff1847} for the graphs of the entry \emph{Verified Enumeration of Trees}~\cite{AFP-Trees},
using the Cauchy--Binet formula of the entry on the Smith normal form~\cite{AFP-SNF}.
\end{abstract}

\paragraph{Acknowledgement.} The formalization and the text of this document were prepared with the
assistance of an AI system, and checked by the author.

\tableofcontents

\section*{Introduction}

Number the vertices other than the root $r$ and the edges, and orient each edge from its smaller to its larger end.
The reduced oriented incidence matrix $N$ has a row for each vertex $v \neq r$ and a column for each edge $e$, with
entry $1$ if $v$ is the smaller end of $e$, $-1$ if it is the larger end, and $0$ otherwise. The reduced Laplacian
$L$ (degrees on the diagonal, $-1$ for each edge) is defined directly, and the proof goes in four steps:
\begin{enumerate}
\item $L = N N^{T}$, entry by entry.
\item By the Cauchy--Binet formula, $\det L$ is the sum of the squares of the maximal minors of $N$, one for each
  set of $|V| - 1$ edges.
\item A minor vanishes if its edges do not connect the graph: some component avoids $r$, and the rows of its
  vertices add up to zero. If its edges form a tree, the square of the minor is $1$, by induction on the tree: a
  leaf contributes a single entry $\pm 1$ to its row or to the column of its edge.
\item A connected set of $|V| - 1$ edges is a spanning tree, and edge sets correspond to sets of column indices.
\end{enumerate}

\parindent 0pt\parskip 0.5ex

\input{session}

\bibliographystyle{abbrv}
\bibliography{root}

\end{document}
